% Generated draft: requires mathematical and source review.
\documentclass[11pt]{article}
\usepackage[margin=25mm]{geometry}
\usepackage{fontspec}
\setmainfont{Latin Modern Roman}
\usepackage{amsmath,amssymb,mathtools,unicode-math}
\setmathfont{Latin Modern Math}
\usepackage{enumitem,needspace,xurl,hyperref}
\setlength{\parindent}{0pt}
\setlength{\parskip}{5pt}
\setlength{\emergencystretch}{3em}
\newcommand{\N}{\mathbb{N}}
\newcommand{\Z}{\mathbb{Z}}
\newcommand{\Q}{\mathbb{Q}}
\newcommand{\R}{\mathbb{R}}
\newcommand{\C}{\mathbb{C}}
\newcommand{\F}{\mathbb{F}}
\newcommand{\A}{\mathbb{A}}
\newcommand{\PP}{\mathbb{P}}
\newcommand{\E}{\mathbb{E}}
\newcommand{\eps}{\varepsilon}
\newcommand{\dd}{\,\mathrm{d}}
\begin{document}
\section*{UnsolvedMath ID: 800369}
\textbf{Source problem code:} catalogue entry
\subsection{Definitions and mathematical statement}
Let \((W,S)\) be a Coxeter system, with length function \(\ell:W\to\Z_{\geq 0}\) and Bruhat order \(\leq\). Thus \(x\leq w\) means that a reduced expression for \(w\) contains a subword that is a reduced expression for \(x\).

Fix an indeterminate \(q^{1/2}\). The Hecke algebra \(\mathcal H(W)\) over \(\Z[q^{1/2},q^{-1/2}]\) has basis \(\{T_w:w\in W\}\), identity \(T_e\), and multiplication determined, for \(s\in S\), by
\[
T_sT_w=\begin{cases}
T_{sw},&\ell(sw)=\ell(w)+1,\\
qT_{sw}+(q-1)T_w,&\ell(sw)=\ell(w)-1.
\end{cases}
\]
Its bar involution is determined by \(\overline{q^{1/2}}=q^{-1/2}\) and \(\overline{T_w}=(T_{w^{-1}})^{-1}\). For \(x\leq w\), the Kazhdan--Lusztig polynomials \(P_{x,w}(q)\in\Z[q]\) are characterized by bar invariance of
\[
C'_w=q^{-\ell(w)/2}\sum_{x\leq w}P_{x,w}(q)T_x,
\qquad \overline{C'_w}=C'_w,
\]
together with \(P_{w,w}(q)=1\) and
\[
\deg P_{x,w}\leq\frac{\ell(w)-\ell(x)-1}{2}\qquad(x<w).
\]
Set \(P_{x,w}(q)=0\) when \(x\not\leq w\).

\textbf{Problem (Brenti, 2024, Problem 7).} Find a combinatorial interpretation for Kazhdan--Lusztig polynomials. In the usual general formulation, this asks for a uniform, explicitly specified combinatorial construction, from \((W,S)\) and \(x\leq w\), of a finite set \(\mathcal A_{x,w}\) and a statistic \(\operatorname{stat}:\mathcal A_{x,w}\to\Z_{\geq0}\) such that
\[
P_{x,w}(q)=\sum_{a\in\mathcal A_{x,w}}q^{\operatorname{stat}(a)}.
\]
Equivalently, the construction would interpret each coefficient \([q^k]P_{x,w}(q)\) as the number of objects of statistic \(k\). The combinatorial objects and their statistic should be specified independently of the numerical coefficients being interpreted.

\textbf{Unresolved scope and interpretation.} The supplied source does not specify a class of Coxeter systems, the permitted combinatorial objects, or a formal criterion for what constitutes a combinatorial interpretation. The all-Coxeter-system counting formulation above makes the customary target explicit; these requirements are an interpretation of the brief source request, not additional conditions quoted from it. The source motivates the problem by reporting a proof of the nonnegativity conjecture and citing \([\mathrm{EW}]\), but the supplied record does not identify that reference bibliographically.
\subsection{Short English statement}
Brenti asks for a combinatorial explanation of the coefficients of Kazhdan--Lusztig polynomials, ideally by constructing objects whose statistic records the exponent of \(q\). The supplied wording leaves the precise scope and admissible kind of interpretation unspecified. Its \texttt{open} label and its reference to a proof of nonnegativity are supplied source assessments, not independently verified conclusions.
\subsection{Sources}
\begin{itemize}
\item Local UnsolvedMath registry, permanent ID \(800369\), exact source problem code \texttt{catalogue entry}, title ``Coxeter Groups and Unimodality --- Problem 7.'' The supplied record attributes the problem to Brenti (2024) and labels it \texttt{open}.
\item Brenti, \emph{Some Open Problems on Coxeter Groups and Unimodality} (2024), Problem 7. Supplied source URL: \url{https://arxiv.org/abs/2410.09897}. The record reports access on 2026-10-16 and extraction by \texttt{arxiv\_latex}; those provenance details are supplied and have not been independently verified here.
\item The supplied statement cites \([\mathrm{EW}]\) for a proof of nonnegativity of Kazhdan--Lusztig polynomials. No authors, title, publication details, or URL for \([\mathrm{EW}]\) are included in the supplied record, so its identity and exact hypotheses remain unverified.
\end{itemize}
\end{document}
